\chapter{Discussion}
\label{chap:discussion}

\section{Analyse critique}

\subsection{Forces}

Le projet atteint l'ensemble de ses objectifs et présente plusieurs qualités
notables. Sur le plan de l'\textbf{ingénierie}, la compilation est propre (aucun
avertissement, aucun conflit grammatical) et la gestion mémoire est exempte de
fuite. Sur le plan de la \textbf{conception}, l'architecture en pipeline avec
front-end commun s'est révélée juste : elle a permis d'ajouter un back-end de
génération de code sans toucher à l'interpréteur, validant a posteriori le choix
d'un AST comme représentation pivot. Sur le plan de la \textbf{robustesse}, toutes
les erreurs prévues sont détectées et localisées, et la campagne de cas limites
n'a révélé aucun plantage non maîtrisé.

\subsection{Faiblesses et limitations}

L'honnêteté impose de recenser les limites, dont certaines sont des choix assumés
et d'autres des bornes de fond de l'approche.

\begin{description}[style=nextline,leftmargin=1em]
  \item[Profondeur de récursion bornée.] L'évaluateur récursant en C, la
        profondeur d'appels \ilang{} est limitée par la pile de l'hôte (rupture
        autour de quelques milliers de niveaux, cf. tableau~\ref{tab:rec}). Une
        machine virtuelle à pile explicite lèverait cette limite.
  \item[Performance de l'interprétation.] Le parcours d'arbre est intrinsèquement
        lent (cf. figure~\ref{fig:bench-fib}). C'est le prix de la lisibilité ; un
        bytecode l'améliorerait sensiblement.
  \item[Absence de portée de bloc.] Les variables déclarées dans un bloc
        (\code{si}, boucle) vivent dans le frame de la fonction, pas dans un frame
        de bloc. C'est cohérent avec la spécification (« globale au principal,
        locale aux fonctions ») mais diffère de langages à portée de bloc.
  \item[Back-end assembleur restreint aux entiers.] Le générateur de code refuse
        les flottants et ne reproduit pas la gestion « douce » de la division par
        zéro (un \code{idiv} par zéro lève un SIGFPE). Choix assumé, hors
        périmètre.
  \item[Pas de système de types.] Le langage est dynamiquement typé au sens
        minimal (entier ou flottant), sans vérification statique.
\end{description}

\subsection{Comparaison avec l'existant}

Par rapport à l'approche antérieure (le mini-langage \textsc{Parlot} et son
analyse manuelle en Python), \ilang{} gagne sur deux tableaux : la grammaire est
formalisée et vérifiable, et l'usage d'outils générateurs élimine des classes
entières d'erreurs (oubli de cas, ambiguïté non détectée). En contrepartie, la
courbe d'apprentissage de Flex et Bison est plus raide qu'une descente récursive
écrite à la main. C'est toutefois un investissement qui se rentabilise dès que la
grammaire se complexifie, et qui paie immédiatement sur le plan de la fiabilité,
puisqu'il supprime des familles entières d'erreurs au lieu de les laisser à la
vigilance du programmeur.

\section{Difficultés rencontrées et solutions}

Cette section revient sur les obstacles techniques marquants et la manière dont
ils ont été levés.

\subsection{Le « dangling else » qui n'existait pas}

Le conflit syntaxique du \emph{dangling else}, c'est-à-dire la question de savoir à
quel \code{si} rattacher un \code{sinon} dans une imbrication, est l'exemple
canonique de conflit shift/reduce enseigné dans tous les cours de compilation.
M'attendant à le rencontrer, j'avais d'abord protégé ma grammaire par une
résolution de priorité, en introduisant un pseudo-token de plus faible priorité
que \code{SINON}. Par acquit de conscience, j'ai ensuite retiré cette protection
pour vérifier si le conflit existait réellement, et Bison n'en a signalé
absolument aucun. Mes deux règles, données ci-dessous, se compilent sans la
moindre ambiguïté.

\begin{lstlisting}[style=gram,caption={Les deux règles du \code{si}, sans aucune directive de priorité.}]
si_alors
    : SI '(' expression ')' bloc                { $$ = new_if($3, $5, NULL); }
    | SI '(' expression ')' bloc SINON bloc     { $$ = new_if($3, $5, $7); }
    ;
\end{lstlisting}

L'explication tient entièrement à un détail de ma grammaire que je n'avais pas
mesuré au départ : les corps de conditionnelle sont \emph{obligatoirement} placés
entre accolades, car la règle \code{bloc} l'impose. Or, dans le cours, le conflit
naît de corps réduits à une instruction nue. Ici, comme tout corps se termine par
une accolade fermante, un \code{sinon} est nécessairement précédé du symbole
\code{\}}. Au moment précis où l'analyseur pourrait hésiter, après avoir reconnu
\code{SI ( expression ) bloc}, le symbole de prévision n'est jamais \code{SINON}
mais toujours \code{\}} ou un début d'instruction. Autrement dit, \code{SINON}
n'appartient pas à l'ensemble \textsc{Follow} de la réduction du \code{si} sans
\code{sinon}, et l'ambiguïté n'a tout simplement pas lieu d'être. J'ai donc
supprimé la directive de priorité devenue inutile. Je retiens de cet épisode une
leçon de méthode : il vaut mieux \emph{tester} une hypothèse avec l'outil que
recopier par réflexe une solution apprise, car le contexte précis de ma grammaire
rendait cette solution superflue.

\subsection{L'ordre d'évaluation et la portée}

La gestion des fonctions a révélé deux subtilités que je n'avais pas anticipées et
qui, toutes deux, touchent à la portée des variables.

La première est une question d'\textbf{ordre}. Lorsqu'on appelle une fonction,
il faut impérativement évaluer les arguments dans l'environnement de l'appelant
\emph{avant} d'empiler le cadre de la fonction appelée. Dans ma toute première
version, j'avais inversé ces deux étapes, et une expression aussi banale que
\code{fibonacci(n - 1)} échouait : une fois le nouveau cadre empilé, la variable
\code{n} était cherchée dans ce cadre encore vide au lieu du cadre de l'appelant.
La correction a consisté à séparer nettement les phases, comme le montre l'extrait
suivant.

\begin{lstlisting}[style=c,caption={L'ordre correct : évaluer les arguments avant d'empiler le cadre.}]
/* 1) arguments evalues dans le cadre de l'APPELANT */
for (i = 0; i < node->nargs; i++) argvals[i] = eval(node->args[i]);
/* 2) seulement maintenant on empile le cadre de l'appele */
env_push();
for (i = 0; i < decl->nargs; i++) env_set(decl->args[i]->name, argvals[i]);
\end{lstlisting}

La seconde subtilité, plus profonde, concerne la \textbf{nature de la portée}. Ma
première implémentation de la recherche de variable parcourait toute la pile
d'appel, du sommet jusqu'à la base. Or ce parcours complet réalise en fait une
portée \emph{dynamique} : une fonction pouvait y lire les variables locales de la
fonction qui l'avait appelée, ce qui n'est pas le comportement attendu d'un langage
moderne. Le test reproducteur ci-dessous affichait \code{123} au lieu d'échouer,
puisque \code{g} voyait le \code{secret} local de \code{f}.

\begin{lstlisting}[style=ilang,caption={Cas qui révèle la portée dynamique (corrigé depuis).}]
fonction g() { retourner(secret) }
fonction f() { secret = 123 retourner(g()) }
afficher(f())          // affichait 123 : g voyait une locale de f !
\end{lstlisting}

La spécification réclamant une portée « lexicale », j'ai corrigé la résolution
pour qu'elle ne consulte que le cadre courant, puis en repli l'environnement
global, sans traverser les cadres intermédiaires de la pile d'appel.

\begin{lstlisting}[style=c,caption={Résolution lexicale corrigée (extrait de \code{symtab.c}).}]
int env_get(const char *name, Value *out) {
    if (lookup_in(current_env, name, out)) return 1;          /* cadre courant */
    if (current_env != global_env && lookup_in(global_env, name, out)) return 1;
    return 0;                                                 /* sinon : absente */
}
\end{lstlisting}

Désormais le programme ci-dessus signale correctement que \code{secret} n'est pas
défini dans \code{g}. Je relève au passage une contradiction du cahier des charges
lui-même, qui qualifiait de « lexicale » une portée dont il décrivait pourtant
l'algorithme dynamique. J'ai tranché en faveur du terme employé, qui correspond
aussi au comportement le plus utile.

\subsection{Le piège du pliage de la division par zéro}

Le pliage de constantes consiste à calculer dès la compilation les
sous-expressions dont les deux opérandes sont des littéraux. Le piège est que
\code{10 / 0} est, lui aussi, une expression à deux opérandes littéraux. Si je le
pliais sans précaution, deux issues fâcheuses étaient possibles : ou bien
l'optimiseur lui-même provoquait une division entière par zéro et un
\code{SIGFPE} au moment de la compilation, ou bien il produisait une valeur
arbitraire qui aurait court-circuité le message d'erreur que l'évaluateur doit
émettre à l'exécution, message exigé par le cahier des charges. La solution
consiste à interdire le pliage des seules opérations qui échoueraient, en laissant
intact le sous-arbre concerné pour que l'évaluateur le rencontre normalement à
l'exécution. C'est le rôle du test \code{can\_fold}, invoqué avant tout pliage :

\begin{lstlisting}[style=c,caption={Le pliage n'est tenté que si l'opération est sûre (extrait de \code{optim.c}).}]
if (node->left->type == NODE_NUM && node->right->type == NODE_NUM
    && can_fold(node->op, node->left->num, node->right->num)) {
    Value v = do_fold(node->op, node->left->num, node->right->num);
    /* ... le noeud BINOP est remplace par le litteral resultat ... */
}
/* sinon : on laisse le noeud tel quel, l'evaluateur tranchera a l'execution */
\end{lstlisting}

Ce petit garde-fou illustre une idée importante : une optimisation ne doit jamais
changer le comportement observable d'un programme, pas même ses messages
d'erreur.

\subsection{Détails d'environnement et de robustesse}

Plusieurs petits obstacles, sans gloire mais instructifs, ont jalonné le travail.
Je les mentionne car leur résolution méthodique a directement contribué à la
propreté finale du logiciel.

Le premier est apparu dès la toute première compilation, sous la forme d'un
avertissement déroutant sur la fonction \code{strdup} :

\begin{lstlisting}[style=gram]
implicit declaration of function 'strdup'
assignment to 'char *' from 'int' makes pointer from integer without a cast
\end{lstlisting}

La cause n'est pas une faute de ma part mais une subtilité de la norme :
\code{strdup} ne fait pas partie du C standard, c'est une extension POSIX. Compilé
en \code{-std=c11} strict, l'en-tête masque son prototype, le compilateur suppose
alors un retour de type \code{int}, et sur une machine 64~bits le pointeur se
trouve tronqué. La correction tient en un mot dans le \code{Makefile} :

\begin{lstlisting}[style=gram]
CFLAGS = -Wall -Wextra -std=gnu11 -g    # gnu11 expose les extensions POSIX
\end{lstlisting}

Le deuxième obstacle, encore plus discret, fut une ligne de dépendance du
\code{Makefile} indentée par mégarde avec une espace, là où Make attend soit une
tabulation pour une recette, soit aucune indentation pour une dépendance.
Le troisième fut le tampon interne de Flex, jamais libéré, que Valgrind signalait
comme « encore atteignable » et que l'appel \code{yylex\_destroy()} en fin de
\code{main} a fait disparaître. Le dernier, enfin, fut un débordement silencieux :
le littéral \code{9999999999} s'affichait \code{1410065407} parce que le lexer le
convertissait avec \code{atoi}, limité à 32~bits, alors que mes valeurs sont des
\code{long}. La correction a consisté à remplacer la conversion :

\begin{lstlisting}[style=gram,caption={Avant et après, dans \code{ilang.l}.}]
{CHIFFRE}+   { yylval.num = atoi(yytext); ... }    /* avant : tronque a 32 bits */
{CHIFFRE}+   { yylval.num = strtol(yytext, NULL, 10); ... }  /* apres : 64 bits */
\end{lstlisting}

Aucun de ces points n'est spectaculaire pris isolément, mais c'est l'attention
portée à chacun qui distingue un programme qui « marche sur les exemples » d'un
programme réellement soigné.

\section{L'extension : du parcours d'arbre au code natif}

Le développement du back-end assembleur, bien que hors périmètre, mérite une
discussion car il éclaire la nature même du projet.

\subsection{Modèle et choix}

Le générateur applique la convention cdecl et matérialise, en assembleur, ce que
l'évaluateur réalisait abstraitement : la pile d'appels devient une vraie pile de
cadres \code{\%ebp}, les paramètres et locales des déplacements, les expressions
une suite d'opérations sur \code{\%eax} et la pile. Les structures de contrôle se
traduisent en labels et sauts ; \code{arrêter} et \code{continuer} deviennent
naturellement des \code{jmp} vers les labels de fin et de continuation de boucle.

\subsection{Obstacles spécifiques}

Trois difficultés propres à la cible native ont dû être traitées : l'installation
de la chaîne 32~bits (\code{gcc-multilib}) ; un avertissement de l'éditeur de
liens (\code{DT\_TEXTREL}) dû à l'adressage absolu incompatible avec le PIE,
résolu en liant avec \code{-no-pie} ; et l'alignement de la pile requis par
\code{printf}, assuré par un \code{andl \$-16, \%esp} en tête de \code{main} et
l'arrondi des cadres à 16~octets.

\subsection{Enseignement}

Voir le \emph{même} AST se transformer tantôt en valeurs (interprétation), tantôt
en instructions machine (compilation), rend tangible toute la théorie de la
chaîne de compilation. C'est, de l'avis de l'auteur, la partie la plus formatrice
du projet, et cela tient sans doute précisément au fait qu'elle dépassait la
commande.

\section{Compromis effectués}

En synthèse, les principaux compromis assumés sont : la \emph{lisibilité} au
détriment de la performance, avec un parcours d'arbre plutôt qu'une machine à
bytecode ; la \emph{simplicité} au détriment de l'exhaustivité, sans portée de
bloc, sans typage statique et avec un back-end limité aux entiers ; et la
\emph{fidélité à la spécification} partout où elle était claire. Là où le cahier
des charges se contredisait lui-même, j'ai fait des choix explicites et
documentés. Pour la portée, j'ai retenu le sens lexical conforme au terme employé
plutôt que l'algorithme dynamique décrit. Pour le point-virgule, présent dans la
BNF mais absent des programmes de démonstration, j'ai préféré le rendre optionnel
afin d'honorer les deux sources à la fois plutôt que d'en sacrifier une. Ces choix
sont cohérents avec la finalité pédagogique du projet.
